\newpage
\begin{center}
\large\textbf{Минобрнауки России}

\large\textbf{Юго-Западный государственный университет}
\vskip 1em
\normalsize{Кафедра программной инженерии}
\vskip 1em
\ifВКР{
        \begin{flushright}
        \begin{tabular}{p{.4\textwidth}}
        \centrow УТВЕРЖДАЮ: \\
        \centrow Заведующий кафедрой \\
        \hrulefill \\
        \setarstrut{\footnotesize}
        \centrow\footnotesize{(подпись, инициалы, фамилия)}\\
        \restorearstrut
        «\underline{\hspace{1cm}}»
        \underline{\hspace{3cm}}
        20\underline{\hspace{1cm}} г.\\
        \end{tabular}
        \end{flushright}
        }\fi
\end{center}
\vspace{1em}
  \begin{center}
  \large
\ifВКР{
ЗАДАНИЕ НА ВЫПУСКНУЮ КВАЛИФИКАЦИОННУЮ РАБОТУ
  ПО ПРОГРАММЕ БАКАЛАВРИАТА}
  \else
ЗАДАНИЕ НА КУРСОВУЮ РАБОТУ (ПРОЕКТ)
\fi
\normalsize
  \end{center}
\vspace{1em}
{\parindent0pt
  Студента \АвторРод, шифр\ \Шифр, группа \Группа
  
1. Тема «\Тема\ \ТемаВтораяСтрока»
\ifВКР{
утверждена приказом ректора ЮЗГУ от \ДатаПриказа\ № \НомерПриказа
}\fi.

2. Срок предоставления работы к защите \СрокПредоставления

3. Исходные данные для создания программной системы:

3.1. Перечень решаемых задач:}

\renewcommand\labelenumi{\theenumi)}

\begin{enumerate}
\item провести анализ предметной области и существующих программных средств построения интерактивных приложений;
\item сформулировать требования к компонентно-ориентированной программной системе и определить порядок ее использования разработчиком;
\item спроектировать архитектуру программной системы, включающую ядро компонентной модели, управление приложением и сценами, графическую подсистему, обработку ввода, скрипты и базовую физику;
\item реализовать основные классы, компоненты, системы обработки и демонстрационные сцены;
\item выполнить модульное, регрессионное и системное тестирование разработанной программной системы.
\end{enumerate}

{\parindent0pt
  3.2. Входные данные и требуемые результаты для программы:}

\begin{enumerate}
\item Входными данными для программной системы являются: исходный код приложения на языке Java, описывающий сцены, сущности, компоненты, системы обработки и пользовательские скрипты; параметры запуска приложения; графические ресурсы, включая шейдеры, текстуры и геометрические данные; события пользовательского ввода с клавиатуры и мыши; параметры объектов сцены, камер, источников света, физических тел и коллайдеров.
\item Требуемыми результатами работы программы являются: сформированное настольное интерактивное приложение на языке Java; отображаемые кадры активной сцены; обновленное состояние сущностей и компонентов сцены; результаты обработки пользовательского ввода; результаты выполнения пользовательских скриптов; результаты базовой обработки физических взаимодействий; диагностические сообщения и отладочная визуализация; демонстрационные сцены, подтверждающие работоспособность основных подсистем.
\end{enumerate}

{\parindent0pt

  4. Содержание работы (по разделам):
  
  4.1. Введение.
  
  4.2. Анализ предметной области: интерактивные приложения как класс программных систем, программные средства построения интерактивных приложений, компонентно-ориентированный подход к организации объектов сцены, разработка интерактивных приложений на Java, анализ существующих решений.
  
4.3. Техническое задание: основание для разработки, цель и назначение разработки, требования к программной системе, требования к архитектурной модели, управлению приложением и сценами, графической подсистеме, пользовательскому вводу, скриптам, базовой физике, демонстрационным сценам, программному интерфейсу, нефункциональным характеристикам и оформлению документации.

4.4. Технический проект: общие сведения о программной системе, обоснование выбора технологий, внутреннее устройство программной системы, структура пакетов, жизненный цикл приложения и сцены, проектирование ядра компонентной модели, графической подсистемы, подсистемы ввода, скриптов, ресурсов, базовой физики и демонстрационных сцен.

4.5. Рабочий проект: спецификация реализованных компонентов и классов программной системы, описание математических, графических, ресурсных и физических классов, систем обработки и демонстрационных сцен, модульное, регрессионное и системное тестирование программной системы.

4.6. Заключение.

4.7. Список использованных источников.

5. Перечень графического материала:

Лист 1. Сведения о ВКРБ

Лист 2. Цель и задачи разработки

Лист 3. Архитектура программной системы

Лист 4. Компонентная модель программной системы

Лист 5. Жизненный цикл кадра приложения

Лист 6. Результаты разработки и тестирования

\vskip 2em
\begin{tabular}{p{6.8cm}C{3.8cm}C{4.8cm}}
Руководитель \ifВКР{ВКР}\else работы (проекта) \fi & \lhrulefill{\fill} & \fillcenter\Руководитель\\
\setarstrut{\footnotesize}
& \footnotesize{(подпись, дата)} & \footnotesize{(инициалы, фамилия)}\\
\restorearstrut
Задание принял к исполнению & \lhrulefill{\fill} & \fillcenter\Автор\\
\setarstrut{\footnotesize}
& \footnotesize{(подпись, дата)} & \footnotesize{(инициалы, фамилия)}\\
\restorearstrut
\end{tabular}
}

\renewcommand\labelenumi{\theenumi.}
